\documentclass[11pt]{article}

\usepackage[margin=1in]{geometry}
\usepackage{amsmath,amssymb}
\usepackage{booktabs}
\usepackage{tabularx}
\usepackage{array}
\usepackage{enumitem}
\usepackage[T1]{fontenc}
\usepackage{hyperref}

\newcommand{\code}[1]{\texttt{#1}}
\newcommand{\Uh}{\hat U}
\newcommand{\tf}{t_{\mathrm{form}}}
\newcolumntype{L}{>{\raggedright\arraybackslash}X}
\newcolumntype{C}{>{\centering\arraybackslash}X}
\emergencystretch=3em

\title{Interacting Brownian particles in \code{src\_deankow/interaction}:\\
formation, merging and splitting of clusters}
\author{FHDeX}
\date{}

\begin{document}
\maketitle

\noindent This note describes the model solved in
\code{src\_deankow/interaction} and gives estimates of four things, for each
interaction potential in the code:
\begin{enumerate}[nosep]
  \item the linear stability of the uniform state and the time for clusters
        to form;
  \item how many clusters form initially;
  \item the time scales for clusters to merge;
  \item the time scales for clusters to split.
\end{enumerate}
The potentials are generalized exponential (\code{gema}), Hegselmann--Krause
(\code{hk}), generalized Morse (\code{morse}) and piecewise parabolic
(\code{pp}). For \code{gema} and \code{hk}, both attractive and repulsive
cases are treated. \code{README.md} lists the inputs and \code{STABILITY.md}
has more on the linear stability and the time step. All numbers are for 2D and
the unit periodic box, with the parameter sets of Section~\ref{sec:params}.

\section{Model}

\subsection{Particles}
$N$ Brownian particles at positions $\mathbf{X}_i$ in the periodic box
interact through a pair potential $U$, weighted by $1/N$:
\begin{equation}
  d\mathbf{X}_i = -\frac1N\sum_j\nabla U(|\mathbf{X}_i-\mathbf{X}_j|)\,dt
     + \sqrt{2D}\,d\mathbf{W}_i .
\end{equation}
The mobility is 1, so $D$ plays the role of temperature. A particle at
$\mathbf{x}$ feels the potential
$C(\mathbf{x}) = N^{-1}\sum_j U(|\mathbf{x}-\mathbf{X}_j|)$.

\subsection{Mean field and Dean--Kawasaki SPDE}
With $\phi = N^{-1}\sum_i\delta(\mathbf{x}-\mathbf{X}_i)$, normalized to
$\int\phi=1$ (so the mean is $\mu_0 = 1/A = 1$),
\begin{equation}
  \partial_t\phi = \nabla\cdot\bigl(D\nabla\phi + \phi\nabla C\bigr)
     + N^{-1/2}\,\nabla\cdot\bigl(\sqrt{2D\phi}\,\boldsymbol{\xi}\bigr),
  \qquad C = U\star\phi .
\end{equation}
This is the McKean--Vlasov equation plus Dean--Kawasaki noise.
\begin{itemize}
  \item Without noise it is the gradient flow of
        $\mathcal F[\phi] = D\!\int\!\phi\ln\phi
        + \tfrac12\!\iint U(\mathbf{x}-\mathbf{y})\phi(\mathbf{x})\phi(\mathbf{y})$.
  \item Its stationary states are fixed points of $\phi = e^{-C/D}/Z$.
  \item $N$ enters only through the noise; the mean-field strength is set by
        $U$, $D$ and $\mu_0$.
\end{itemize}
In the code, $C$ is computed by FFT (\code{AdvancePhiAtLevel.cpp}), and the
flux $D\nabla\phi+\phi\nabla C$ is discretized on faces (\code{mykernel.H}).
The particle version (\code{StochasticPC}) sums the pair forces within the
cutoff \code{amr.ip\_range}.

\subsection{Potentials}
\begin{align}
  \text{\code{gema}:}\quad & U = \varepsilon\,e^{-(r/R)^\alpha}, \\
  \text{\code{hk}:}\quad & U = \tfrac12\varepsilon R\,(1-r^2/R^2) \ \ (r<R),\quad 0 \ \ (r\ge R), \\
  \text{\code{morse}:}\quad & U = -\varepsilon_{\mathrm{att}}\,e^{-(r-r_e)/R_{\mathrm{att}}}
      + \varepsilon_{\mathrm{rep}}\,e^{-(r-r_e)/R_{\mathrm{rep}}}, \\
  \text{\code{pp}:}\quad & U = -\varepsilon R\,w(r/R),\quad
      w(s) = \begin{cases}
        \alpha(s^2-a^2)/2 + \beta(a^2-1)/2, & s\le a,\\
        \beta(s^2-1)/2, & a<s\le1,\\
        0, & s>1 .
      \end{cases}
\end{align}
\begin{center}
\begin{tabular}{@{}lll@{}}
\toprule
\code{ip\_type} & range & attractive when \\
\midrule
\code{gema} & unbounded, fast decay & $\varepsilon<0$ \\
\code{hk} & $R$ & $\varepsilon<0$ \\
\code{morse} & exponential tails & always ($\varepsilon_{\mathrm{att}}>\varepsilon_{\mathrm{rep}}$, $R_{\mathrm{att}}>R_{\mathrm{rep}}$ required) \\
\code{pp} & $R$ & $\varepsilon<0$ \\
\bottomrule
\end{tabular}
\end{center}
For \code{hk} and \code{pp}, $\varepsilon=-\gamma$ and $R=\ell$ map onto the
potentials $\gamma\ell\,w(x/\ell)$ of Gerber et al.\ (arXiv:2510.17629).

\section{General results}

\subsection{Linear stability and the formation time}
Linearizing about $\phi=\mu_0$ gives, for each Fourier mode $\mathbf k$,
\begin{equation}
  \sigma(k) = -k^2\bigl(D + \mu_0\,\Uh(k)\bigr), \qquad
  \Uh(k) = 2\pi\int_0^\infty U(r)\,J_0(kr)\,r\,dr .
\end{equation}
The uniform state is unstable at $k$ when $\mu_0\Uh(k) < -D$. On the periodic
box only the lattice modes $\mathbf k = 2\pi(n_x,n_y)$ exist; $k=0$ is the
conserved mass.
\begin{itemize}
  \item \textbf{Attractive kernels} ($\Uh(0)<0$): the band reaches down to the
        smallest box mode. Growth peaks where $k^2(-D-\mu_0\Uh(k))$ is
        largest, a finite $k$ set by the kernel's shape. The system clusters
        and then coarsens.
  \item \textbf{Repulsive kernels} ($\Uh(0)>0$): only kernels whose $\Uh$ has
        a negative lobe can go unstable, at the lobe's $k^*$, when
        $D<D_{\mathrm{crit}}=\mu_0|\Uh_{\min}|$. The result is a cluster
        crystal.
\end{itemize}
The noise seeds each mode with relative amplitude about $N^{-1/2}$. Nonlinear
clusters appear once the fastest mode reaches $O(1)$:
\begin{equation}
  \tf \approx \frac{\ln\sqrt N}{\sigma_{\max}} = \frac{\ln N}{2\,\sigma_{\max}} ,
\end{equation}
up to an $O(1)/\sigma_{\max}$ correction.

\subsection{Number of clusters}
The fastest mode sets the spacing $\lambda = 2\pi/|\mathbf k_{\max}|$, and the
pattern breaks into blobs on a lattice:
\begin{equation}
  n_{\mathrm{square}} = \Bigl(\frac{|\mathbf k_{\max}|L}{2\pi}\Bigr)^2, \qquad
  n_{\mathrm{hex}} = \frac{\sqrt3}{8\pi^2}\,A\,|\mathbf k_{\max}|^2 \approx 0.87\,n_{\mathrm{square}} .
\end{equation}
A hexagonal lattice with reciprocal vectors of magnitude $|\mathbf k_{\max}|$
has nearest-neighbour spacing $a = 4\pi/(\sqrt3|\mathbf k_{\max}|)$. In small
boxes a specific $(n_x,n_y)$ mode wins, so $n_{\mathrm{square}}$ is usually
closer. The count is uncertain by about $\pm20\%$ when several modes grow at
nearly the same rate.

The cluster mass is $m\approx1/n$. The width follows from the self-consistent
well $C(r)\approx m\,U(r)$:
\begin{itemize}
  \item harmonic core: $\sigma_c^2 = D/(mU''(0))$;
  \item flat core (\code{gema}, $\alpha>2$, $U\approx U(0)+|\varepsilon|(r/R)^\alpha$):
        $r_c\approx R\,(D/m|\varepsilon|)^{1/\alpha}$.
\end{itemize}

\subsection{Representative parameter sets}\label{sec:params}
\begin{center}
\begin{tabularx}{\textwidth}{@{}l L l l L@{}}
\toprule
case & parameters & $D$ & $N$ & source \\
\midrule
gema attr. & $\varepsilon=-666$, $R=0.1$, $\alpha=3$ & 0.5 & 81920 & \code{inputs\_fv}, sign flipped \\
gema rep. & $\varepsilon=+666$, $R=0.1$, $\alpha=3$ & 0.25 ($D_{\mathrm{crit}}=0.55$) & 81920 & \code{inputs\_fv}, lower $D$ \\
hk attr. & $\varepsilon=-3000$, $R=0.1$ & 1 & 81920 & \code{inputs\_fv\_hk\_cluster} \\
hk rep. & $\varepsilon=+2\cdot10^4$, $R=0.1$ & 0.5 ($D_{\mathrm{crit}}=0.92$) & 81920 & --- \\
morse & $\varepsilon_{\mathrm{att}}=200$, $\varepsilon_{\mathrm{rep}}=100$, $R_{\mathrm{att}}=0.1$, $R_{\mathrm{rep}}=0.05$ & 0.5 & 81920 & --- \\
pp & $\varepsilon=-7300$, $R=0.45$, $a=0.27$, $\alpha=1$, $\beta=0.02$ & 1 & 327680 & \code{inputs\_fv\_pp\_merge} \\
\bottomrule
\end{tabularx}
\end{center}

\subsection{Merging mechanisms}
\paragraph{(M1) Drift.} Clusters within each other's range attract. Each moves
toward the other at speed $m|U'(d)|$, so a pair at distance $d_0$ merges in
\begin{equation}
  t_{\mathrm{drift}} \approx \int_{\text{core}}^{d_0}\frac{dd}{2m|U'(d)|},
\end{equation}
independent of $N$. This is the fast route when $\lambda$ lies inside the range.

\paragraph{(M2) Brownian coalescence.} Clusters beyond range diffuse with
$D/(Nm)$ and merge when random motion closes the gap $g\le\lambda$:
\begin{equation}
  t_{\mathrm{brown}} \approx \frac{g^2 N m}{8D} \qquad (\propto N).
\end{equation}

\paragraph{(M3) Mass exchange (Ostwald-type ripening).} Particles escape over
the well depth $m|U(0)|$. Smaller clusters, with shallower wells, lose mass
faster and eventually vanish:
\begin{equation}
  t_{\mathrm{exch}} \approx \frac{\sigma_c^2}{D}\,e^{m|U(0)|/D} .
\end{equation}
This is independent of $N$ and present in the mean-field PDE: the dynamical
metastability of Gerber et al.

\subsection{Splitting}
\begin{itemize}
  \item \textbf{Attractive kernels.} Separating the two halves of a cluster
        costs a total energy $N(m/2)^2|U(0)|$, so the rate is about
        $e^{-Nm^2|U(0)|/4D}$. That is exponentially small in $N$, and absent
        from the gradient-flow PDE. What looks like splitting during
        formation is the break-up of stripes into blobs, on a time of about
        $1/\sigma_{\max}$.
  \item \textbf{Repulsive cluster crystals.} The number of clusters is fixed
        by the spacing $2\pi/k^*$, and the occupancy by $N/n$. Occupancies
        change by single-particle hops over the lattice barrier $\Delta C$, at
        a rate of about $(D/a^2)e^{-\Delta C/D}$. Creating or removing a site
        is a collective, activated rearrangement, much slower still.
\end{itemize}

\section{Results for each potential}

\subsection{Summary}
\begin{center}
\small
\begin{tabularx}{\textwidth}{@{}l c c c c L c c c@{}}
\toprule
case & $\sigma_{\max}$ & mode & $\tf$ & $n$ (sq./hex) & spacing & $m$ & width & $m|U(0)|/D$ \\
\midrule
gema attr. & 3633 & (3,2) & 0.0016 & 13 / 11 & 0.277 ($2.8R$) & 0.077 & $r_c\approx0.021$ & 102 \\
gema rep. & 757 & (7,4) & 0.0075 & 65 / 56 & hex $a=0.143$ ($1.43R$) & 0.015 & 0.008 & --- \\
hk attr. & 265 & (3,1) & 0.021 & 10 / 9 & 0.316 ($3.2R$) & 0.10 & 0.018 & 15 \\
hk rep. & 1750 & (10,3) & 0.0032 & 109 / 94 & hex $a=0.111$ ($1.11R$) & 0.009 & ring, $r\approx0.012$ & --- \\
morse & 283 & (1,1) & 0.020 & 2 & 0.707 ($7.1R_{\mathrm{att}}$) & 0.5 & 0.007 & 100 \\
pp & 299 & (3,1) & 0.021 & 10 / 9 & 0.316 ($0.70R$) & 0.10 & 0.025 & 15 \\
\bottomrule
\end{tabularx}

\medskip
\begin{tabularx}{\textwidth}{@{}l L L L L@{}}
\toprule
case & drift merge (M1) & Brownian (M2) & mass exchange (M3) & splitting \\
\midrule
gema attr. & none: force $\approx7\cdot10^{-6}$ & $\approx120$ & $\propto e^{102}$: never & $e^{-1.6\cdot10^5}$: never \\
gema rep. & clusters repel & --- & hops: $\Delta C/D=17$, $\tau_{\mathrm{hop}}\approx3\cdot10^6$ & lattice change: never \\
hk attr. & none ($\lambda>R$) & $\approx100$ & $\approx3\cdot10^{-4}e^{15}\approx10^3$ & $e^{-3\cdot10^4}$: never \\
hk rep. & clusters repel & --- & hops: $\Delta C/D=12$, $\tau_{\mathrm{hop}}\approx4\cdot10^3$ & $\gg\tau_{\mathrm{hop}}$ \\
morse & $\approx0.06$ & $\approx5\cdot10^3$ & $\propto e^{100}$: never & $e^{-10^6}$: never \\
pp & $\approx0.015$ & $\approx400$ & $\approx2\cdot10^3$ & $e^{-1.2\cdot10^5}$: never \\
\bottomrule
\end{tabularx}
\end{center}
At $128^2$, $\Delta t\approx2$--$4\cdot10^{-6}$ for these runs, so $t=100$ is
about $3\cdot10^7$ steps. Only times up to about 0.1--1 are routinely
reachable.

\subsection{\texorpdfstring{GEM-$\alpha$, attractive ($\varepsilon<0$)}{GEM-alpha, attractive}}
\begin{itemize}
  \item \textbf{Stability.} $\Uh(0)=0.0284\,\varepsilon$ for $\alpha=3$, so
        $\Lambda=38$. The band runs from the smallest box mode to about
        $6\cdot2\pi$.
  \item \textbf{Formation.} Very fast: $\tf\approx0.0016$. Growth peaks at
        $kR\approx2.3$ ($\lambda\approx0.28$), independent of $|\varepsilon|$
        once $\Lambda\gg1$.
  \item \textbf{Clusters.} About 11--13, scaling like $1/R^2$. The core is flat,
        so clusters are compact blobs with steep edges and
        $r_c\approx R(D/m|\varepsilon|)^{1/3}\approx0.02$.
  \item \textbf{Merging.} The force at $2.8R$ is negligible
        ($e^{-(2.8)^3}\approx3\cdot10^{-10}$), and mass exchange is ruled out
        by $m|\varepsilon|/D\approx100$. Only Brownian coalescence remains,
        with $t\approx120$ at $N=81920$, growing like $N$. The pattern is
        effectively frozen on simulation times; weaker attraction or higher
        $D$ moves $\lambda$ toward $R$, where drift merging becomes possible.
  \item \textbf{Splitting.} Never.
\end{itemize}

\subsection{\texorpdfstring{GEM-$\alpha$, repulsive ($\varepsilon>0$)}{GEM-alpha, repulsive}}
\begin{itemize}
  \item \textbf{Stability.} For $\alpha>2$, $\Uh$ has a negative lobe: its
        minimum is $-0.029\,\Uh(0)$ at $kR=5.03$. The uniform state is
        unstable for $D<D_{\mathrm{crit}}=0.55$ at $\varepsilon=666$.
  \item \textbf{Formation.} At $D=0.25$, $\sigma_{\max}\approx760$ and
        $\tf\approx0.0075$. Near onset $\sigma\propto(D_{\mathrm{crit}}-D)k^{*2}$,
        so $\tf$ diverges.
  \item \textbf{Clusters.} A hexagonal cluster crystal with $a\approx1.43R$:
        about 56 clusters of about 1300 particles. The number is set by $R$
        alone ($k^*R\approx5$); the occupancy is $N/n$.
  \item \textbf{Merging and splitting.} Neighbouring clusters repel. Particles
        hop over $\Delta C\approx4.4$ ($\Delta C/D\approx17$, $\tau_{\mathrm{hop}}\approx3\cdot10^6$),
        so the lattice is static. Defects from the formation stage persist.
\end{itemize}

\subsection{\texorpdfstring{HK, attractive ($\varepsilon<0$)}{HK, attractive}}
\begin{itemize}
  \item \textbf{Stability.} $\Uh(k)=\varepsilon R^3\hat S(kR)$ with
        $\hat S(0)=\pi/4$; $\Lambda=2.28$. The onset is
        $|\varepsilon|\approx4D/(\pi R^3)$ for small $R$.
  \item \textbf{Formation.} The fastest mode is $(3,1)$, with $\sigma\approx265$
        and $\tf\approx0.021$.
  \item \textbf{Clusters.} About 10, with $\lambda\approx3.2R$; in the
        strong-coupling limit $\lambda\to2.06R$. Either way the spacing is
        outside the range, and the compact support freezes the pattern.
        Gaussian clusters with $\sigma_c^2=DR/(m|\varepsilon|)$, giving
        $\sigma_c\approx0.018$.
  \item \textbf{Merging.} No drift. Brownian coalescence takes $\approx100$
        ($\propto N$). Mass exchange, with $m|U(0)|/D=m\gamma R/2D\approx15$,
        takes $\approx10^3$; this is the $e^{\gamma\ell\Delta m}$
        metastability of Gerber et al. Both are out of reach of simulations.
  \item \textbf{Splitting.} Never.
\end{itemize}

\subsection{\texorpdfstring{HK, repulsive ($\varepsilon>0$)}{HK, repulsive}}
\begin{itemize}
  \item \textbf{Stability.} The lobe $\hat S_{\min}=-0.046$ at $kR=6.38$ gives
        $D_{\mathrm{crit}}=0.92$ at $\varepsilon=2\cdot10^4$.
  \item \textbf{Formation.} At $D=0.5$, $\sigma_{\max}\approx1750$ and
        $\tf\approx0.003$.
  \item \textbf{Clusters.} A hexagonal crystal with $a\approx1.11R$: about 94
        clusters of about 870 particles. Each parabolic piece is concave
        (curvature $-\varepsilon/R$), so the lattice potential's minimum lies
        off the site, at about 0.012. The clusters are small rings, confined
        by the kernel edges, which need $\Delta x\ll0.01$ to resolve.
  \item \textbf{Merging and splitting.} Hops cross $\Delta C/D\approx12$,
        with $\tau_{\mathrm{hop}}\approx4\cdot10^3$; the crystal is static.
\end{itemize}

\subsection{Morse}
\begin{itemize}
  \item \textbf{Stability.} $\Uh$ increases wherever it is negative, so there
        is either long-wave demixing ($\Uh(0)<0$) or stability; no crystal.
        The code defaults ($\varepsilon_{\mathrm{att}}=1$,
        $\varepsilon_{\mathrm{rep}}=0.5$) give $\Uh(0)=-0.055$, stable for
        $D\gtrsim0.06$. The set here gives $\Uh(0)=-10.5$ and $\Lambda\approx21$.
  \item \textbf{Formation.} With $\Uh\approx\Uh(0)+ck^2$, the fastest mode
        satisfies $k_{\max}^2\approx(\mu_0|\Uh(0)|-D)/(2c)$: here $(1,1)$, with
        $\sigma\approx283$ and $\tf\approx0.02$.
  \item \textbf{Clusters.} About 2. More need a larger $|\Uh(0)|/D$ or a shorter
        $R_{\mathrm{att}}$; near threshold, one band. Here
        $\varepsilon_{\mathrm{att}}/R_{\mathrm{att}}=\varepsilon_{\mathrm{rep}}/R_{\mathrm{rep}}$,
        so the core is harmonic with
        $U''(0)=-\varepsilon_{\mathrm{att}}/R_{\mathrm{att}}^2+\varepsilon_{\mathrm{rep}}/R_{\mathrm{rep}}^2$
        and $\sigma_c\approx0.007$. Otherwise the core has a cusp and
        $r_c\approx D/(m|U'(0)|)$.
  \item \textbf{Merging.} The exponential tail reaches between clusters, so
        drift merging takes $\approx0.06$ and the system coarsens quickly
        toward one cluster.
  \item \textbf{Splitting.} Never.
\end{itemize}

\subsection{PP (piecewise parabolic)}
\begin{itemize}
  \item \textbf{Stability.} $\Uh$ is the sum of two HK transforms, a tail of
        radius $R$ and a core of radius $aR$. The core sets the fastest mode
        when the tail is weak ($\beta\lesssim0.02$ here).
  \item \textbf{Formation.} The fastest mode is $(3,1)$, with
        $\sigma\approx299$ and $\tf\approx0.021$. Observed: clusters from
        $t\approx0.014$, ten by $t\approx0.018$.
  \item \textbf{Clusters.} About 10 predicted, and 10 observed. The spacing,
        $0.32=0.70R$, lies inside the tail.
        $\sigma_c^2=DR/(m|\varepsilon|\alpha)$, giving $\sigma_c\approx0.025$.
  \item \textbf{Merging.} The tail pulls neighbours together with
        $m|U'(\lambda)|\approx10$, so drift merging takes $\approx0.015$,
        independent of $N$. Observed: nine mergers, from 10 clusters to 1,
        between $t\approx0.027$ and $0.036$.
  \item \textbf{Splitting.} Never.
\end{itemize}

\section{Scaling summary}
\begin{center}
\begin{tabularx}{\textwidth}{@{}l L@{}}
\toprule
quantity & scaling \\
\midrule
$\tf$ & $\ln N/(2\sigma_{\max})$, with $\sigma_{\max}\sim(\mu_0|\Uh|-D)k^2$ \\
number of clusters & $\propto k_{\max}^2A$: about $A/R^2$ times a kernel constant; nearly independent of $N$, and of $\varepsilon$ far above onset \\
drift merging & $\propto1/(m|U'(\lambda)|)$; independent of $N$; needs $\lambda$ inside the range \\
Brownian merging & $\propto Nm\lambda^2/D$ \\
mass exchange & $\propto e^{m|U(0)|/D}$ \\
splitting (attractive) & $\propto e^{Nm^2|U(0)|/4D}$: never \\
hopping (repulsive) & $\propto(a^2/D)\,e^{\Delta C/D}$ \\
\bottomrule
\end{tabularx}
\end{center}
\paragraph{How the estimates were made.}
\begin{itemize}
  \item $\sigma(k)$ is the numerical Hankel transform of $U$ on the box modes.
        These are continuum values; the code's startup printout gives the
        discrete ones, within a few percent.
  \item Mass, width and drift use $m\approx1/n$ and a harmonic or flat well.
  \item The mass-exchange and hopping times are Arrhenius estimates with an
        $O(1)$ prefactor. Only their exponents are reliable.
  \item The lattice barriers are from direct sums over a hexagonal lattice of
        point clusters.
\end{itemize}

\begin{thebibliography}{9}
\bibitem{dean} D.~S.~Dean, \emph{J. Phys. A} \textbf{29}, L613 (1996).
\bibitem{likos} C.~N.~Likos, A.~Lang, M.~Watzlawek and H.~Löwen,
  \emph{Phys. Rev. E} \textbf{63}, 031206 (2001).
\bibitem{gerber} N.~Gerber, R.~S.~Gvalani, M.~Hairer, G.~A.~Pavliotis and
  A.~Schlichting, Formation of clusters and coarsening in weakly interacting
  diffusions, arXiv:2510.17629.
\end{thebibliography}

\end{document}
